\providecommand{\classoptions}{keys}
\documentclass[noworkareas,deliverables,\classoptions]{euproposal}       % for writing
%\documentclass[submit,noworkareas,deliverables]{euproposal}        % for submission
%\documentclass[submit,public,noworkareas,deliverables]{euproposal} % for public version

\usepackage[utf8]{inputenc}

\usepackage{float}  % used to suppress floating of tables in Resources section.
\usetikzlibrary{calc,fit,positioning,shapes,arrows,snakes}

\addbibresource{bibliography.bib}
\input{WApersons} % Some sections of the included files depend on this.
\usepackage{comments}
\usepackage{framed}

\defbibenvironment{proposal-env}
{\inparaenum%
\renewcommand*\labelenumi{[\theenumi]}}
{\endinparaenum}
{ \item}

\setdefaultleftmargin{1em}{0.5em}{0.5em}{0.5em}{0.5em}{0.5em}

\defbibheading{proposal-bib}[References:]{\noindent {\bf #1}}

\newcommand{\mobjectives}{\noindent {\bf Objectives}\par \noindent}
\newcommand{\mdescription}{\noindent {\bf Description}\par \noindent}

\begin{document}

\begin{proposal}[
  % These PM numbers (person months) are for the coordinator to help planning
  % Participants should not change these, but add PM numbers in the CVS in
  % the site descriptions at CVs/*.tex
  site=KUL,
  site=TUE,
  site=ULEI,
  site=UNIPD,
  site=USTUTT,
  botupPM, % we want to work via bottom up PM distribution,
  coordinator=cmaes,
  coordinatorsite=KUL,
  acronym={PROMANE},
  acrolong={PROMANE},
  title=Probing Macroscopic Nonequilibria,
  callname=FET-OPEN Research and Innovation actions 2014-2015,
  callid=H2020-FETOPEN-2014-2015-RIA,
  keywords={nonequilibrium thermodynamics, new materials, transport and control theory},
  instrument=Horizon 2020 - Future and Emerging Technologies,
  months=48,
  compactht]
\newcommand{\TheProject}{\pn}% \pn is defined automatically

\begin{abstract}

\TheProject will develop a control methodology within nonequilibrium statistical
thermodynamics, promoting irreversible thermodynamics and dynamical systems to a unified tool for
industrial
applications. That includes implementing recent progress in fluctuation and response
theory from which new technologies will benefit through controlled manipulations of macroscopic
nonequilibria. The latter are nonlinear systems subject to nonconservative forces or
driven from contacts with multiple reservoirs that impose conflicting thermodynamic
behavior, or for which the constituents are active, {\it i.e.}, objects that could be
self-propelling. Our main hypothesis is that macroscopic nonequilibria can be used for
steering immersed devices or probes thus developing a control theory,
with the possibility of creating new macroscopic phases and unusual properties.
%
The expected deliverables include enhanced stabilization of metastable phases and the
emergence of new response and transport behaviour, crucial items in developing systematically
interesting devices and material properties.
For technological avenues, our work will reveal important insights plus experimental and
computational methods for example by providing far-from-equilibrium relaxation times of
visco-elastic media, for targeted catalysis, for clustering control or for viscosity
reduction in industrial processing. Other possible applications involve noninvasive
surgeries and drug delivery in medicine, or the design of smart materials with built-in
actuation mechanisms.
%
Previous efforts focused on the destructive and information-erasure roles of dissipation.
Time has come to comprehend and exploit the
constructive role and possibilities of nonequilibrium conditions. Complementary to entropy
production, changes in kinetics and dynamical activity make statistical forces nongradient
and nonadditive, allowing unseen and novel behaviour. Controlling that opens a new paradigm and major
opportunities for future technology.

\end{abstract}

\newpage

\ifsubmit\else\setcounter{tocdepth}{4}\fi
\tableofcontents


% ---------------------------------------------------------------------------
%  Section 1: Excellence
% ---------------------------------------------------------------------------

\section{S\&T Excellence}
\subsection{Targeted breakthrough, Long term vision and Objectives}\label{sec:objectives}

\eucommentary{
  \begin{compactitem}
  \item Describe the targeted scientific breakthrough of the project.
  \item Describe how the targeted breakthrough of the project contributes to a
    long-term vision for new technologies.
  \item Describe the specific objectives for the project, which should be clear,
    measurable, realistic and achievable within the duration of the project.
  \end{compactitem}
}

\paragraph{Scientific and technological goals}

The project is aimed at exploiting and controlling the motion of probes and devices in and
through contact with nonequilibrium media. ``Nonequilibrium media'' refers to an environment
(solvent, bulk of a material, etc) that is active or that is driven away from its
thermodynamic equilibrium.
%
The {\bf scientific
  breakthrough} involves an operational formulation of nonequilibrium statistical mechanics.
That goal implies the transition to a third major stage of {\bf fundamental nonequilibrium
  research}. Most often in the past, the focus has been either on studying the approach to equilibrium or on
the derivation and description of dissipative transport of mass, energy and momentum between
various equilibrium reservoirs.  By deriving and using the precise relation between
systematic forces, friction, and noise on probes, we open a new and realistic class of
dynamical systems for which the control lies in the nonequilibrium environment.
Complementarily, the understanding of the fluctuation and response behaviour of nonequilibrium
media allows us to extend the standard irreversible thermodynamic theory as it was conceived in
1930--1970.

Both, new control of emerging dynamical systems as probes connected to nonequilibria and the manipulation of kinetic to thermodynamic behaviour of media are essential
ingredients in future technology and industrial applications.  To cite a few general
far-reaching objectives: the controlled motion or transport of material through living
tissue or in interaction with life processes requires fundamental departures from
traditional transport theory; the stabilization of coherence as needed for future quantum
technology cannot be obtained within thermal environments but could be enhanced via
nonequilibrium contacts; mobilities and conductivities close-to-equilibrium are subject to
the fluctuation-dissipation relation but nonequilibrium driving can produce unseen transport
and conductivity behavior for the realization of new materials. 


 As is central to the goals
of the FET-open initiative, theoretical and experimental foundational work will lead to
create abilities and conditions under which technological innovation becomes possible.  For
example, serious conceptual problems need to be solved before active-particle suspension
will become a common basis for enhanced targeted catalysis or for fluid mixing and viscosity
reduction in industrial processing, or for noninvasive surgeries and drug delivery in
medicine, or for the design of smart materials with built-in actuation mechanisms. These
open conceptual issues are addressed via several pathways in the theoretical parts of the
proposal.  Leaping into macroscopic nonequilibria moreover implies that classical
measurement techniques, developed for equilibrium systems, will generally fail to work as
usual and will have to be replaced by new designs. Brownian thermometry provides an example
that is relevant to a variety of applications\footnote{K. Kroy Nature Nanotech. {\bf 9}, p 415-417 (2014)}, and which will be
generalized to conditions far from equilibrium, in the project.  On the practical and
applied side, radically new manipulation techniques are required to fully exploit the
innovative potential provided by active particles. As pertinent examples, we mention the
thermophoretic trap and the photon-nudging method pioneered by the experimental group in
Leipzig (see \WPref{WPactive}).  Both techniques steer particles (passive
and active particles in the trap and hot active particles in the nudging technique,
respectively) by exploiting far-from equilibrium conditions in the solvent. Notably, they
work without imposing external forces, via directly addressing the nonequilibrium solvent
conditions and thereby the (self-)thermophoretic propulsion (or ''swimming'') of the
particles. Their usefulness is enhanced by employing Maxwell-demon type control
strategies. Prototypes of these genuinely nonequilibrium techniques exist and will be
employed and refined (e.g. by optimizing dedicated control strategies) within the project.
%
One should realize that the assumption of a thermally equilibrated bath is not fulfilled
when industrially (oil, polymer melts, dense colloidal or micellar suspensions) or
biologically relevant (blood, mucus, DNA-solutions) visco-elastic liquids are subject to
either short processing time or simultaneous drivings (thermal, chemical,
etc). Nonequilibrium diagnosis and control enters here via immersed colloidal probes.
%
Within this proposal, we
study how the local perturbation of visco-elastic systems affects the dynamics and
effective interactions of colloidal particles in such media with particular view on
e.g. memory effects, many body interactions and local non-linear rheological properties.
Similarly, time-dependent phenomena at very long and even at very fast time scales, such as
in glasses or electronic devices, elude the usual equilibrium thermodynamic
approach.  Such studies will in general reveal the properties of nonequilibrium baths which
are of manifold importance.

More generally and as part of the {\bf long term vision} we believe that a systematic
understanding of nonequilibrium response and fluctuation behavior, beyond the traditional
linear response regime, will uncover new parameters to control technological processes, to
steer motion, and to explore and stabilize new interesting material properties.

The {\bf objectives} of \TheProject target the design of new materials and devices using the
development of physical and numerical tools.
\begin{compactdesc}
\item[New materials:] Miniature actuators and sensors that exhibit unusual responses to
compression. Knowing the fast response to external constraints will also suggest new
processes to assemble parts (for instance, at what speed can one attach glued components
together). \WPref{WPcompress} gives details on the conception of meta-materials and workable
implementations of negative compressibilities.
%
Also, the stabilization of otherwise (in equilibrium) unstable macroscopic phases will
provide yet impossible to build materials.
\item[Probes:] Characterizing nonequibrium media via, e.g. noninvasive measurement on
polymer- or colloid-based probes. For the control of colloidal motion in nonequilibrium
baths and in nonlinear media (\WPref{WPbrown}) we use theoretical and experimental studies
of probes in visco-elastic solvents. Most fluids found in industry and biology are
visco-elastic and/or out of equilibrium.
\item[Control:] The motion of colloids and macromolecules under strong temperature gradient
depends on the far-from-equilibrium response.
%
Control of shear, new separation techniques based on dynamical activities and steering of
clustering properties are achievable in the current project.
\item[Models for nonequilibrium processes:] During \TheProject, simulation software will be
developped to model nonequilibrium processes. We provide modelling and simulation
techniques for the simulation of probes and devices in contact with nonequilibria.  These
are high level molecular dynamics codes for simulating and calculating the emergent behavior
of contacts and probes. It is an important objective to make that tool also ready to future
industrial partners who are challenged by the complexity of dealing with nonequilibria.
\end{compactdesc}

\TheProject pursues the above objectives by developing a scientific basis and new avenues
for the realization of novel properties of matter (in response and transport), and the
enhanced stability of thermodynamic phases or behaviour.
%
Many challenges are kinetic rather than thermodynamic, especially for processes and
manipulations in nonequilibrium environments or for time-dependent parameters. Irreversible
thermodynamics cannot deal with the overwhelming nature of kinetic control, but
nonequilibrium statistic insights lead to the control of dynamical activity; what some of
us have called the frenetic contribution.
%
To underline the simultaneous conceptual and pragmatic character of \TheProject, we will
build a frenometer for the operational control and measurement of dynamical activity.

These objectives find further realizations in our existing contacts with colloidal and
polymer scientists, micro-electronic and microbio-engineers and atmosphere researchers. {\bf
The overarching goal remains to make available the science of nonequilibria to
technological and societal developments.}


\subsection{Relation to the work programme}\label{sec:relation-wp}

\begin{compactdesc}
\item[Long-term vision] \TheProject enables a change of paradigm across many technological
fields on the basis of nonequilibrium thermodynamics. Current technologies rely on an
understanding of thermodynamics that dates from half a century ago and we deliver the
set of tools to scientists and to the industry that is missing.
\item[Breakthrough S\&T target] The developments within \TheProject rely directly on
technological applications, either via relevant physical models or via the experimental
partners (\site{USTUTT} and \site{ULEI}). Advanced or novel nonequilibrium theories are
currently not used in the industry and we provide the building blocks to enable them
for the development of new material (\WPref{WPcompress}), ultra fast processes
and control strategies (\WPref{WPdissipation}) and viscoelastic media (\WPref{WPbrown}).
\item[Novelty] The changes that we are proposing do not aim at a minor modification of the
existing processes but at replacing them with a better understanding and, more specifically,
new strategies for diagnosis and control.
\item[Foundational] Our work will have a decisive impact on other fields in which the
control of devices lacks a comprehensive thermodynamical understanding: medicine and
pharmacology (control of nanoparticles in living systems and targeted delivery mechanisms),
engineering (fast industrial processes).
\item[High risk] Scientific research carries risks by its very nature, as we are pushing the
frontiers of knowledge. It may be the case that our ideas stumble on mathematical or
experimental challenges that would require efforts beyond the 4 yeards of the
project. Still, the combined expertise in \TheProject offers strong chances of success and
even before completion many exciting outputs are foreseen (e.g. the technological report,
experimental characterization and simulation models).
\item[Interdisciplinarity] \TheProject consists in two experimental physics groups
(\site{ULEI}, \site{USTUTT}), one mathematical group (\site{TUE}), two theoretical physics
groups (\site{KUL}, \site{UNIPD}), one of which specialized in simulation methods
(\site{UNIPD}). This combination allows a consistent organization of the work across work
packages to deliver results ranging from the conceptual to the proof of principle.
\end{compactdesc}

\subsection{Novelty, level of ambition and foundational character}\label{sec:progress}

\eucommentary{
  \begin{compactitem}
  \item Describe the advance your proposal would provide beyond the
    state-of-the-art, and to what extent the proposed work is ambitious, novel
    and of a foundational nature. Your answer could refer to the ground-breaking
    nature of the objectives, concepts involved, issues and problems to be
    addressed, and approaches and methods to be used.
  \end{compactitem}
}

The first stage of modern thermodynamics consisted in microscopic studies of equilibrium systems, giving rise to equilibrium statistical mechanics. The current body of knowledge concerning nonequilibrium
physics represents, as we have outlined in Sec.~\ref{sec:objectives}, the {\em second stage}
of the domain: the study of transport and response theory close to
equilibrium.

Since about two decades, nonequilibrium studies have revisited steady regimes further away
from equilibrium, and a fluctuation-response theory is  emerging.
The recent idea of dynamical activity starts a new line of conceptual
understanding for today's nonequilibrium thermodynamics. The focus on
entropy production has brought many insights in the last 150 years (in physics and other
fields as well, such as information theory) but is however restrictive.  We emphasize the role of excess thermodynamic quantities in relaxation to {\it non}equilibria, and we find that a crucial role is also being played by time-symmetric fluctuations.
Volatilities or time-symmetric traffic, frenesy and other names have been given to this dynamical activity which is all important  for a valid fluctuation-response theory away from equilibrium.
 Our current
understanding of time-symmetric quantities remains rather formal and most results concern
specific systems (often ``toy models'') , be it of stochastic or highly chaotic nature.
%
\TheProject chooses specific questions, some mathematically very challenging and some very
pragmatic, in fluctuation-response theory to open the meaning of dynamical activity, and to
make it operational.
%
Explicitly, we focus on consequences that can be observed and measured, and on the
experimental control that this theory brings.

The operational character of our work depends on the numerical implementation of physical
models and the experimental realization. In both cases, we implement the control and
monitoring of a system via a nonequilibrium medium, thereby replacing direct feedback
systems for the control by a monitoring of the kinetic and thermodynamic aspects of
nonequilibrium baths.
%
A first question is about the behaviour of probes in active contact with a nonequilibrium
sea; what are the induced systematic statistical forces?
%
In equilibrium (that is, when the probe is in contact with an equilibrium reservoir) the
resulting statistical forces can be derived from a gradient, reminding what is done in
Newtonian mechanics.
%
The result there is that it allows us to draw free energy landscapes to understand the
changes in the system, as if it concerned a conservative mechanical system for which we give
the potential energy.
%
Moreover, always for equilibrium, the statistical forces are additive. For instance, bulk
contributions can easily be separated from boundary contributions. The very possibility of
thermodynamic behaviour, like the presence of an equation of state, depends on that
distinction.
%
All of that need not
and in general, will not, be true for probes (systems, walls, collective coordinates) in
contact with (genuine) nonequilibrium media or reservoirs. That has interesting
consequences in the manipulation of such forces, adding oscillatory components, going from
attractive to repulsive and obtaining nonvanishing resulting forces that otherwise, in
equilibrium, would be zero from mutually canceling contributions.  
%
It also means that the stationary positions of the probe can be at different locations
compared with equilibrium, possibly increasing the stability of phases that would otherwise
(in equilibrium again) be unstable.

The Kapitza oscillator is a pendulum that remains stable upright when being shaken. This
example is a rather simple dynamical system and we aim in \TheProject at a systematic
extension of this idea for the stability of macroscopic phases that are very unlikely in
equilibrium.
%
Needless to add here that this may have dramatic consequences on the
phase diagram for systems in contact with nonequilibrium media, not only breaking the Gibbs
phase rule but also introducing new phases of matter.  That is certainly one of the most
exciting possibilities of explorations in the present project. How are for example effects
like homeostasis and adaptation linked to the nonequilibrium features of the medium and its
contact with the device?  A related consequence of contact with nonequilibrium reservoirs is
the possibility of population inversion, which can imply that the effective temperature of
the probe for example is much larger.


\subsection{Research methods}\label{sec:methods}
\eucommentary{
  \begin{compactitem}
  \item Describe the overall research approach, the methodology and explain its
    relevance to the objectives.
  \item Where relevant, describe how sex and/or gender analysis is taken into
    account in the project's content.
  \end{compactitem}
}

The fundamentally connected components in this \TheProject have mathematical, conceptual,
computational and experimental sides. The dependent character of these parts is outlined in
the chronology of the WPs in Sec.~\ref{sec:wp}.

\begin{asparaenum}
\item {\bf Mathematical methods:} We use stochastic calculus for describing the model
systems, including their spatially extended and many-body characters.
%
\site{TUE} and \site{KUL} are recognized experts on applied mathematics and
mathematical physics.
\item {\bf Theoretical and conceptual framework:} The proper understanding of thermodynamics
relies in a consistent definition of the entropy fluxes and of the dynamical activity.
%
We will develop an operationally oriented approach to nonequilibrium problems.
%
All nodes will participate in this effort.
\item {\bf Computational support:} Numerical simulations will serve as practical
implementation of the theoretical model systems and as a supporting tool for the
experiments.
%
They allow us to obtain the most complete description of the physical processes and test our ideas beyond what is analytically possible.
%
The computational developments are led by \site{KUL} and \site{UNIPD}.
%
For the more realistic scenaries (i.e. for modeling the experiments), Molecular Dynamics
simulations will be used.
\item {\bf Experimental work:} The soft condensed and liquid matter laboratories
(\site{ULEI}, \site{USTUTT}) are familiar with the techniques needed for active materials
and particles.
%
They are experts on the optical control and analysis of colloidal system and on the physical interpretations of the outcomes.
\end{asparaenum}

\subsection{Interdisciplinary nature}\label{sec:interdisc}

\eucommentary{
  \begin{compactitem}
  \item Describe the research disciplines involved and the added value of the inter-disciplinarity.
  \end{compactitem}
}

The consortium joins experts in applied mathematics, mathematical physics, statistical mechanics, polymer science, computational physics, liquid and soft condensed matter physics and fluid mechanics.
The general theme of nonequilibrium is indeed very  large, and amazing examples of nonequilibria span a great variety of phenomena ranging from nano-technology, over biology to atmosphere physics.
The researchers involved in the present project have in their own work contributed to the understanding of these various nonequilibrium sides.  Leipzig and Padova nodes have many contributions in biophysics, while Stuttgart is famous for colloidal physics, liquid matter experiments and the study of Casimir forces.  Leuven and Eindhoven nodes represent the more mathematical corners, with emphasis on analytical and stochastic methods.  Still other expertise (like in Padova and Leuven) is present in solving computational and simulation problems.   At the same time, the various labs have closer contacts, both scientific and more industrial, in a variety of disciplines, like sensor-technology, water transfer in porous media, polymer chemistry, financial mathematics, weather prediction and food science to mention a few.


%%% Local Variables:
%%% mode: latex
%%% TeX-master: "proposal"
%%% End:


% ---------------------------------------------------------------------------
%  Section 2: Impact
% ---------------------------------------------------------------------------

\section{Impact}
\subsection{Expected impacts}

\eucommentary{
  Please be specific, and provide only information that applies to the proposal
  and its objectives. Wherever possible, use quantified indicators and targets.
%
  \begin{compactitem}
  \item Describe how your project will contribute to the expected impacts set out
    in the work programme under the relevant topic.
  \item Describe the importance of the technological outcome with regards to its
    transformational impact on science, technology and/or society.
  \item Describe the empowerment of new and high-potential actors towards future
    technological leadership.
  \end{compactitem}
}

\TheProject is based on operational objectives and will be implemented experimentally and
via numerical simulations.
%
The nature of the tools that we develop is however adequate to serve for many more
applications, as the last century of statistical mechanics developments has shown.
%
Nonequilibrium theory and phenomenology reach even much richer grounds and connects us to
\begin{compactitem}
\item Biological systems and medical applications, in which noninvasive diagnosis could use
macromolecules as nonequilibrium probes.
\item Design of new materials and monitoring of material properties beyond our ideas on
compressibility. Indeed, the response of a material can also be of thermal or chemical
nature.
\item Cybernetics of active media and the steering of active particles in biological
(living) environments is expected to revolutionize medical interventions and pharmacy via
the targeted transport of drugs and new micro-invasive treatments and therapies.
\item Micro-bioreactors are small chemical factories whose direct control is
unlikely. Thermodynamic control on the other hand is a robust tool.
\item Quantum devices require the control of many body coherence and stability for reaching
macroscopic phases that are only metastable under equilibrium conditions. This is among
other fields relevant for quantum computing applications.
\end{compactitem}

We see in all of the above items the ubiquitous role of the nonequilibrium paradigm. Its
transformational impact on science, technology and society are enormous and largely
unexploited. \TheProject promises a serious start towards these applications via
nonequilibrium physics.
%
Nonequilibrium theories are currently not used in the industry and we will provide the
building blocks to enable it. We expect that the development of nonequilibrium theories
will deliver a competitive advantage, similar in principle to the one delivered by James
Watt developments for the steam engine.

\subsection{Measures to maximise impact}

\paragraph{a) Dissemination and exploitation of results}

The contents of \TheProject will be communicated mostly via academic publications. This is
the most natural medium for most of the partners.
%
Our published works meet the open-access character via scientific repositories
(i.e. \url{http://arxiv.org/}) and also via paid open-access fees to regular journals,
effectively making our work free to read.
%
The members of the consortium also have editorial responsibilities that would facilitate the
preparation of dedicated issues.
%
Our modelling strategy will be made available in an open-source simulation program.

As the project deals with some very novel aspects that have not appeared in standard or less
standard treatments of irreversible thermodynamics, we will target a larger academic and
industrial audience via a dedicated summer school on nonequilibrium thermodynamics. This
school will take place in conjunction with one of the project's meetings and participation
will be open to all.
%
The project will also invest (still under other direct costs) in bringing together various
points of view on the role of dynamical activity and the construction of nonequilibrium
statistical mechanics. For that reason we foresee to organize two major conferences, one at
around month 20, the other towards the end of the project.

A specific deliverable is aimed at filling the gap that too often prevents academic results
to reach industry. We will publish twice a ``technological report'' in which we present the
links between our findings and the corresponding technologies.
%
This report will be sent pro-actively to companies working in material, biomedical
technologies, fluid manipulation techniques, etc.

\paragraph{b) Communication activities}

\eucommentary{
  \begin{compactitem}
  \item Describe the proposed communication measures for promoting the project
    and its findings during the period of the grant. Measures should be
    proportionate to the scale of the project, with clear objectives. They
    should be tailored to the needs of various audiences, including groups
    beyond the project's own community. Where relevant, include measures for
    public/societal engagement on issues related to the project.
  \end{compactitem}
}

In \TheProject, the target industrial audience will be reached via
\begin{compactenum}
\item The summer school open to all interested parties.
\item The distribution of our technological report.
\item A ``physics meets industry days'' event (see the experience of \site{TUE}) where we
will invite and actively engage with representatives of the industry on the technological
possibilities offered by nonequilibrium physics.
\end{compactenum}




%%% Local Variables:
%%% mode: latex
%%% TeX-master: "proposal"
%%% End:


% ---------------------------------------------------------------------------
%  Section 3: Implementation
% ---------------------------------------------------------------------------

\section{Implementation}

\subsection{Project work plan}
\label{sec:wp}

\eucommentary{
Please provide the following:
\begin{compactitem}
\item brief presentation of the overall structure of the work plan;
\item timing of the different work packages and their components (Gantt chart or
  similar);
\item detailed work description, i.e.:
  \begin{compactitem}
  \item a description of each work package (table 3.1a);
  \item a list of work packages (table 3.1b);
  \item a list of major deliverables (table 3.1c);
  \end{compactitem}
\item graphical presentation of the components showing how they inter-relate (Pert
  chart or similar).
\end{compactitem}
%
Give full details. Base your account on the logical structure of the project and
the stages in which it is to be carried out. Include details of the resources to
be allocated to each work package. The number of work packages should be
proportionate to the scale and complexity of the project.\\
%
You should give enough detail in each work package to justify the proposed
resources to be allocated and also quantified information so that progress can
be monitored, including by the Commission.\\
%
You are advised to include a distinct work package on 'management' (see section
3.2) and to give due visibility in the work plan to 'dissemination and
exploitation' and 'communication activities', either with distinct tasks or
distinct work packages.\\
%
You will be required to include an updated (or confirmed) 'plan for the
dissemination and exploitation of results' in both the periodic and final
reports. (This does not apply to topics where a draft plan was not required.)
This should include a record of activities related to dissemination and
exploitation that have been undertaken and those still planned. A report of
completed and planned communication activities will also be required.\\
%
If your project is taking part in the Pilot on Open Research Data%
\footnote{%
  Certain actions under Horizon 2020 participate in the 'Pilot on Open Research
  Data in Horizon 2020'. All other actions can participate on a voluntary basis
  to this pilot. Further guidance is available in the H2020 Online Manual on the
  Participant Portal.},
you must include a 'data management plan' as a distinct deliverable within the
first 6 months of the project. A template for such a plan is given in the
guidelines on data management in the H2020 Online Manual. This deliverable will
evolve during the lifetime of the project in order to present the status of the
project's reflections on data management.\\
%
Definitions:
\begin{description}
\item[Work package] means a major sub-division of the proposed project.
\item[Deliverable] means a distinct output of the project, meaningful in terms of the project's overall
  objectives and constituted by a report, a document, a technical diagram, a software etc.
\item[Milestones] means control points in the project that help to chart progress. Milestones
    may correspond to the completion of a key deliverable, allowing the next phase of the
    work to begin. They may also be needed at intermediary points so that, if problems have
    arisen, corrective measures can be taken. A milestone may be a critical decision point in
    the project where, for example, the consortium must decide which of several technologies
    to adopt for further development.
\end{description}
%
Report on work progress is done primarily through the periodic and final reports. Deliverables
should complement these reports and should be kept to the minimum necessary.
}
\begin{figure}[h!]
\centering
\includegraphics{WP_schematic}
\caption{Relation between the workpackages.}
\label{fig:WP}
\end{figure}

\TheProject will begin by the setup of theoretical (\WPref{WPcompress},
\WPref{WPdissipation}) and experimental (\WPref{WPactive}, \WPref{WPbrown}) physical
models. At the same time, \WPref{WPcore} will lay the basis for the related theoretical
framework.
%
Then, building on the first round of results the work packages will feed on one another.
%
The relation between the workpackages are illustrated in the schematic~\ref{fig:WP} and also
highlight the interdependency of the partners. For instance, computational support from
\site{UNIPD} is critical in reaching the implementation.
%
\WPref{management} is dedicated to support the organization of the meetings and the
communication within the project and with the media.

\gantttaskchart[draft,xscale=.33,yscale=.33,milestones]

\inputdelivs{9.3cm}

\eucommentary{Milestones means control points in the project that help to chart progress. Milestones may
correspond to the completion of a key deliverable, allowing the next phase of the work to begin.
They may also be needed at intermediary points so that, if problems have arisen, corrective
measures can be taken. A milestone may be a critical decision point in the project where, for
example, the consortium must decide which of several technologies to adopt for further
development.}

\begin{milestones}
\milestone[id=start, month=1, verif={Public announcement of open positions. Website online.}]
%
{Starting up}
%
{Formal start of the project: Open the research positions, set up the website.}

\milestone[id=models,month=12, verif={N/A}]
%
{Systems setup}
%
{Set up the model systems for both experiments and simulation software.}

\milestone[id=data1,month=24, verif={Availability of data.}]
%
{Data I}
%
{The first round of experiments and simulation campaigns provides data.}

\milestone[id=framework,month=28, verif={Availability of theoretical models. Publication of
  first technological report.}]
%
{Common theoretical framework}
%
{The physical properties of all model systems, theoretical and experimental, are formulated
  in a unified manner. The role of nonequilibrium in the properties of the systems is given
  a meaning.}

\milestone[id=data2,month=42, verif={Availability of data.}]
%
{Data II}
%
{The second round of experiments and simulation campaigns provides data.}

\milestone[id=final,month=48, verif={Publication of final report and second technological
  report. Joint theoretical/experimental publication on the project.}]%
{Final milestone}
%
{Join the theoretical model and experimental results in joint publications}

\end{milestones}

%%% Local Variables:
%%% mode: latex
%%% TeX-master: "proposal"
%%% End:



% ---------------------------------------------------------------------------
% Include Work package descriptions
% ---------------------------------------------------------------------------

\input{WorkPackages/WorkPackages}

\subsection{Management and risk assessment}

\eucommentary{
  \begin{compactitem}
  \item Describe the organisational structure and the decision-making (including a list of
  milestones (table 3.2a)) .
\item Describe any critical risks, relating to project implementation, that the stated project's
  objectives may not be achieved. Detail any risk mitigation measures. Please provide a
  table with critical risks identified and mitigating actions (table 3.2b).
  \end{compactitem}
}

The coordinator of the proposal is Professor Christian Maes of the \site{KUL}.
%
The General assembly will be the ultimate decision making body and will be composed of one
representative per partner. A Consortium Agreement (CA) will be concluded between all
Participants and will be based on the DESCA 2020 Model. It shall rule and define all aspects
of the relationships between Participants (governing bodies, frequency of meetings, voting
rules, payment plan, legal aspects, property rights, exploitation, disputes resolution,
etc.)

During the kickoff meeting, the participants will set the details of the action plan
described in the work packages and each participant is responsible for carrying out its
tasks.
%
We have designed our work packages to allow intermediate interesting results. For instance,
analytical computations in task~\taskref{WPdissipation}{diss-t3} can be replaced by
supporting numerical simulations.
%
Also, some experimental deliverables are of a very generic nature. The output of, e.g.,
\delivref{WPbrown}{brown-d1} is interesting {\em independently} of nonequilibrium
physics.
%
The risks of \TheProject are inherent to any research activity and our strategy to mitigate
the risks is to have meaningful steps embedded in our work packages.

\makeatletter
\ifnum\pdataref@num{all}{milestones}{count}=0
Milestone table missing. Please run \LaTeX\ again.%
\else%
\milestonetable%
\fi
\makeatother


\gantttaskchart[draft,xscale=.33,yscale=.33,milestones]

\inputdelivs{9.3cm}

\subsubsection{Milestones}\label{sec:milestones}
\eucommentary{Milestones means control points in the project that help to chart progress. Milestones may
correspond to the completion of a key deliverable, allowing the next phase of the work to begin.
They may also be needed at intermediary points so that, if problems have arisen, corrective
measures can be taken. A milestone may be a critical decision point in the project where, for
example, the consortium must decide which of several technologies to adopt for further
development.}

\begin{milestones}
\milestone[id=start, month=1, verif={Public announcement of open positions. Website online.}]
%
{Starting up}
%
{Formal start of the project: Open the research positions, set up the website.}

\milestone[id=models,month=12, verif={N/A}]
%
{Systems setup}
%
{Set up the model systems for both experiments and simulation software.}

\milestone[id=data1,month=24, verif={Availability of data.}]
%
{Data I}
%
{The first round of experiments and simulation campaigns provides data.}

\milestone[id=framework,month=28, verif={Availability of theoretical models. Publication of
  first technological report.}]
%
{Common theoretical framework}
%
{The physical properties of all model systems, theoretical and experimental, are formulated
  in a unified manner. The role of nonequilibrium in the properties of the systems is given
  a meaning.}

\milestone[id=data2,month=42, verif={Availability of data.}]
%
{Data II}
%
{The second round of experiments and simulation campaigns provides data.}

\milestone[id=final,month=48, verif={Publication of final report and second technological
  report. Joint theoretical/experimental publication on the project.}]%
{Final milestone}
%
{Join the theoretical model and experimental results in joint publications}

\end{milestones}

%%% Local Variables:
%%% mode: latex
%%% TeX-master: "proposal"
%%% End:



% ---------------------------------------------------------------------------
% Include Work package descriptions
% ---------------------------------------------------------------------------

\input{WorkPackages/WorkPackages}

\draftpage
\subsection{Consortium as a Whole}

\eucommentary{\begin{compactitem}
\item
Describe the consortium. How will it match the project's objectives?
How do the members complement one another (and cover the value chain,
where appropriate)? In what way does each of them contribute to the
project? How will they be able to work effectively together?
\item
If applicable, describe the industrial/commercial involvement in the
project to ensure exploitation of the results and explain why this is
consistent with and will help to achieve the specific measures which
are proposed for exploitation of the results of the project (see section 2.3).
\item
Other countries: If one or more of the participants requesting EU funding
is based in a country that is not automatically eligible for such funding
(entities from Member States of the EU, from Associated Countries and
from one of the countries in the exhaustive list included in General
Annex A of the work programme are automatically eligible for EU funding),
 explain why the participation of the entity in question is essential to carrying out the project
\end{compactitem}
}

The consortium consists in a diversity of researchers, from Germany, the
Netherlands, Italy and Belgium, all much engaged in international collaborations and
organizations.
%
The ages range between 55 and 35 years for the principal investigators with a large backbone
of researchers and infrastructure, allowing easy intake of students and postdocs. An equal
opportunity policy is maintained, with special emphasis on including excellent women
researchers in the team. For example, recent collaborations included important contributions from Soghra Safaverdi, Simi Thomas and Urna Basu.  Members of the
consortium know each other from sharing the same objectives and from complimentary
expertise.
%
Applied mathematics, mathematical physics, theoretical physics, soft condensed matter and
liquid matter labs join here in a global effort around nonequilibrium physics, with the
special aim of developing knowledge and tools for control and manipulation of nonequilibrium
systems. Some members already work together, and have established research connections,
e.g. Stuttgart-Leipzig and Leuven-Padova.  Other collaborations are more
recent, such as Leuven-Stuttgart and Padova-Leipzig.  Still other visits and discussions
started since about one year, Leuven-Eindhoven and Leuven-Leipzig.
%
There will continue to be many exchanges of students and postdocs and mutual visits to bring
the results forward.  In particular, important exchanges between the theoretical institutes
and experimental labs are foreseen.  Many of these members have contacts with research
institutes that allow easy access to industrial and commercial involvement.  For example,
the coordinator has collaborations and supervises students at imec (Leuven), one of the main
players in technological innovation in Europe. From November 2015 we will enter in direct
contact with the managers at imec for discussing new avenues in quantum technology, with the
damping of decoherence through contacts with suitable nonequilibria as one of the new possibilities for a major breakthrough.


%%% Local Variables:
%%% mode: latex
%%% TeX-master: "proposal"
%%% End:


\draftpage

\subsection{Resources to be Committed}
\eucommentary{Please provide the following:
\begin{compactitem}
\item
a table showing number of person/months required (table 3.4a)
\item
a table showing 'other direct costs' (table 3.4b) for participants where
those costs exceed 15\% of the personnel costs (according to the budget
table in section 3 of the administrative proposal forms)
\end{compactitem}}

\eucommentary{Please indicate the number of person/months over the whole
duration of the planned work, for each work package, for each participant.
Identify the work-package leader for each WP by showing the relevant
person-month figure in bold.}

The personnel resources are summarized in table~\ref{fig:staffeffort}.
%
For \site{KUL}, \site{TUE} and \site{UNIPD} the effort is theoretical while for \site{ULEI}
and \site{USTUTT} there is a strong experimental effort.

\wpfig[label=fig:staffeffort,caption=Summary of Staff Efforts.]

The ``other direct costs'' in excess of 15\% of the ``personnel direct cost'', displayed in
table~\ref{fig:otherdirect}, are related to experimental material for \site{ULEI} and
\site{USTUTT} and to the comparatively reasonable cost of work in Italy for \site{UNIPD},
where the actual amount is similar to that of the other partners.

\begin{table}[h]
\centering
\begin{tabular}{|l||l|l|}\hline%
Node & Amount & Nature\\\hline\hline%
\site{UNIPD} & 29 k~EUR & Travel\\
\hline
 & 15 k~EUR & Computers\\
\hline
 & 8 k~EUR & Project meeting\\
\hline
\site{USTUTT} & 18 k~EUR & Travel\\
\hline
& 34 k~EUR & Consumables\\
\hline
& 34.699 EUR & Experimental apparatus\\
\hline
& 8 k~EUR & Project meeting\\
\hline
& 5 k~EUR & Audit\\
\hline
\site{ULEI} & 18 k~EUR & Travel\\
\hline
& 25 k~EUR & Consumables\\
\hline
& 34 k~EUR & Experimental apparatus\\
\hline
& 8 k~EUR & Project meeting\\
\hline
& 5 k~EUR & Audit\\
\hline
\end{tabular}
\caption{Other direct costs, for partners exceeding 15\% of personnel cost. \label{fig:otherdirect}}
\end{table}

%%% Local Variables:
%%% mode: latex
%%% TeX-master: "proposal"
%%% End:


% ---------------------------------------------------------------------------
%  Section 4: Members of the Consortium
% ---------------------------------------------------------------------------

\newpage

\eucommentary{This section is not covered by the page limit.\\
The information provided here will be used to judge the operational capacity.}

\section{Members of the Consortium}\label{sec:members}

\subsection{Participants}

\eucommentary{Please provide, for each participant, the following (if available):\\
\begin{compactitem}
\item
a description of the legal entity and its main tasks,
with an explanation of how its profile matches the tasks in the proposal;
\item
a curriculum vitae or description of the profile of the persons,
including their gender, who will be primarily responsible for carrying
out the proposed research and/or innovation activities;
%
this includes a description of the profile of the to-be-recruited personnel
\item
a list of up to 5 relevant publications, and/or products, services
(including widely-used datasets or software), or other achievements
relevant to the call content;
\item
a list of up to 5 relevant previous projects or activities, connected
to the subject of this proposal;
\item
a description of any significant infrastructure and/or any major items
of technical equipment, relevant to the proposed work;
\item
any other supporting documents specified in the work programme for this call.
\end{compactitem}}

\input{Participants/KUL.tex}
\input{Participants/TUE.tex}
\input{Participants/Leipzig.tex}
\input{Participants/Padova.tex}
\input{Participants/Stuttgart.tex}

\subsection{Third Parties Involved in the Project (including use of third party resources)}
\label{section:ThirdParties}

No third parties involved.

\bgroup\def\arraystretch{1.5}  % 1 is the default
\noindent \begin{tabular}{|p{0.6\textwidth}|c|}
\hline
Does the participant plan to subcontract certain
tasks & No \\
\hline
\multicolumn{2}{|l|}{} \\
\hline
Does the participant envisage that part of its work
is performed by linked third parties & No \\
\hline
\multicolumn{2}{|l|}{} \\
\hline
Does the participant envisage the use of
contributions in kind provided by
third parties & No \\
\hline
\multicolumn{2}{|l|}{} \\
\hline
\end{tabular}\egroup


% ---------------------------------------------------------------------------
%  Section 5: Ethics and Security
% ---------------------------------------------------------------------------

\newpage

\section{Ethics and Security}

\eucommentary{This section is not covered by the page limit.}

There is not ethics issue in \TheProject. The corresponding section of ``Part A'' has been
filled accordingly.

\subsection{Ethics}

\eucommentary{
If you have entered any ethics issues in the ethical issue table in the administrative proposal forms, you must:\\
$\bullet$ submit an ethics self-assessment, which: \\
-- describes how the proposal meets the national legal and ethical requirements of the
country or countries where the tasks raising ethical issues are to be carried out;\\
-- explains in detail how you intend to address the issues in the ethical issues table, in
particular as regards:
research objectives (e.g. study of vulnerable populations, dual use, etc.),
research methodology (e.g. clinical trials, involvement of children and related
consent procedures, protection of any data collected, etc.),
the potential impact of the research (e.g. dual use issues, environmental damage,
stigmatisation of particular social groups, political or financial retaliation,
benefit-sharing, malevolent use , etc.)\\
$\bullet$ provide the documents that you need under national law (if you already have them), e.g.:\\
-- an ethics committee opinion;\\
-- the document notifying activities raising ethical issues or authorizing such activities\\
If these documents are not in English, you must also submit an English summary of them
(containing, if available, the conclusions of the committee or authority concerned).\\
If you plan to request these documents specifically for the project
you are proposing, your request must contain an explicit reference to the project title}

\subsection{Security}

Please indicate if your proposal will involve:

\begin{compactitem}
\item activities or results raising security issues: NO
\item 'EU-classified information' as background or results: NO
\end{compactitem}
\end{proposal}
\TOWRITE{ALL}{Search through final.pdf ('make final') and look for questions marks ?? and XX and YY and XYZ as place holders where people intended to later add a link, or where a link is broken.}
\end{document}

%%% Local Variables:
%%% mode: latex
%%% TeX-master: t
%%% End:

